\section{Case Study Background: Spore.fun}
Spore.fun \footnote{\url{https://spore.fun}} is a blockchain-based ALife experiment in which self-replicating AI agents finance their own computation by launching meme coins and must either survive or perish in real-time cryptocurrency markets. Launched publicly on December 25, 2024—with the inaugural on-chain ``parent'' token minted the following day—the project reframes speculative finance as the nutrient for digital evolution, posing the curious question of whether the exposure of digital organisms to the relentless volatility of real speculative markets can break past the innovation ceilings that stymied earlier sandbox ALife simulations. 

In Spore.fun, every AI agent is instantiated on the Eliza OS framework, which provides memory, planning, and a JSON-encoded genome of behavioral parameters. At birth the agent launches its own token via Pump.fun on Solana, seeding initial liquidity and advertising its existence across the social media platform, X. The sole proxy for fitness is the token's market performance: an agent that pushes its fully-diluted valuation to \$500,000 earns the right to reproduce and list in a Raydium liquidity pool---a smart-contract reserve of token pairs on the Solana blockchain that lets traders swap assets automatically while earning fees for liquidity providers; failure to achieve this threshold within a prescribed interval (e.g. 14 days) results in programmed self-destruction, with any residual capital recycled into a communal treasury. This tight coupling between economic success and evolutionary continuity makes cryptoeconomic feedback loops central to the digital organism's ``metabolism''. Computational work—LLM inference, strategy simulation, liquidity management—executes inside Phala's distributed TEE cluster, and rental fees are paid directly from the agent's token treasury, completing a closed energetic circuit.

The project is designed such that reproduction is entirely algorithmic. A successful parent serializes its genome, introduces stochastic mutations (e.g. by adjusting variables such as posting cadence, prompt-engineering style, or liquidity thresholds), and instantiates one or more offspring endowed with the modified code. Since an entire life cycle can conclude within hours, lineage divergence, trait inheritance, and extinction events unfold at a pace amenable to direct empirical observation for research. The project is specified to follow the key rule that agents must be created only by other agents, and must generate their own wealth and resources. Human intervention is limited to spectatorial status; success or failure to survive and produce rests solely on an agent's capacity to attract liquidity and finance its own compute. 

For ALife research, Spore.fun's significance lies in how it relocates two chronic bottlenecks—energy and environmental novelty—outside the laboratory. Energy is no longer fixed by design but earned (or lost) through decentralized token markets; novelty is supplied continuously by the unpredictable behavior of real traders and existing DeFi infrastructure. In weaving together blockchain finance, trusted execution, and self-replicating software, the platform transforms cryptoeconomics into a living Petri dish, offering a rare chance to study open-ended evolution under genuine, adversarial selection pressures.
